\section*{Capítulo \ref{ch:b2:fluid-kinematics} --- Cinemática dos fluidos}

\begin{solution}{exo:b2:fluid-kinematics:1}
(a) $\ell/L \sim 7 \times 10^{-8}$: contínuo. (b) $0.07$: no limite --- as
leis do contínuo começam a falhar (deslizamento nas paredes). (c) $\ell = 70 \times 10^{-9}
/3 \times 10^{-7} = \qty{0.23}{m}$, $\ell/L = 0.23$: não há contínuo; o gás é
uma chuva de moléculas independentes (regime molecular livre).
\end{solution}

\begin{solution}{exo:b2:fluid-kinematics:2}
$a = v\,\dd v/\dd x = v_0^2(1 + x/L)/L$ (de $v_0^2/L$ a $2v_0^2/L$).
$\dd x/\dd t = v_0(1 + x/L)$: $t = (L/v_0)\ln2 = 0.69\,L/v_0$, entre
$L/2v_0$ e $L/v_0$.
\end{solution}

\begin{solution}{exo:b2:fluid-kinematics:3}
$D_V = \int_0^Rv_{\max}(1 - r^2/R^2)\,2\pi r\,\dd r = \tfrac12\pi R^2v_{\max}$;
$\langle v\rangle = v_{\max}/2 = \qty{0.20}{m/s}$; $D_V = \qty{1.6e-5}{m^3/s} =
\qty{16}{mL/s}$.
\end{solution}

\begin{solution}{exo:b2:fluid-kinematics:4}
(a) $\operatorname{div} = 0$, $\vect\omega = \vect 0$: incompressível,
irrotacional. (b) $\operatorname{div} = 0$, $\omega_z = 2\Omega$: rotação
sólida. (c) $\operatorname{div} = 2a \ne 0$, irrotacional (uma fonte
plana). (d) $\operatorname{div} = 0$, $\omega_z = -k$: cisalhamento. (e)
$\operatorname{div} = b$: incompressível se e só se $b = 0$; irrotacional para
todo $b$.
\end{solution}

\begin{solution}{exo:b2:fluid-kinematics:5}
$D_V = \qty{0.20}{L/s}$. Mangueira, $S = \qty{1.13}{cm^2}$: $v = \qty{1.8}{m/s}$; bico,
$S = \qty{7.1}{mm^2}$: $v = \qty{28}{m/s}$. Aceleração média $\Delta(v^2/2)/L
= (28^2 - 1.8^2)/0.08 \approx \qty{1e4}{m/s^2} \approx 1000\,g$. Tempo $\approx
\int\dd x/v$: alguns milissegundos (\qty{4}{ms} para um perfil linear de celeridade).
\end{solution}

\begin{solution}{exo:b2:fluid-kinematics:6}
(a) $\dd y/\dd x = (V/U)\cos\omega t$: retas de coeficiente angular
$(V/U)\cos\omega t$, todas paralelas, que se inclinam com o tempo. (b) $x = Ut$, $y =
(V/\omega)\sin\omega t$: a senoide $y = (V/\omega)\sin(\omega x/U)$. (c) Em
$t = 0$ as \omterm{def:b2:fluid-kinematics:euler}{linhas de corrente} estão a \qty{45}{\degree} e a \omterm{def:b2:fluid-kinematics:euler}{trajetória} deixa
a origem a \qty{45}{\degree}; em $t = \pi/2\omega$ as \omterm{def:b2:fluid-kinematics:euler}{linhas de corrente} são
horizontais, ao passo que a \omterm{def:b2:fluid-kinematics:euler}{trajetória} continua a mesma senoide. Elas diferem porque
o escoamento não é permanente: uma \omterm{def:b2:fluid-kinematics:euler}{linha de corrente} é um instantâneo, uma \omterm{def:b2:fluid-kinematics:euler}{trajetória} é uma história.
(d) $\vect a = \partial_t\vect v = (0, -V\omega\sin\omega t)$ --- o campo é
uniforme, de modo que o termo convectivo se anula.
\end{solution}

\begin{solution}{exo:b2:fluid-kinematics:7}
(a) $D_V = 50 \times 2 \times 1.2 = \qty{120}{m^3/s}$. (b) $120/60 = \qty{2.0}{m/s}$.
(c) $150/150 = \qty{1.0}{m/s}$. (d) $\rho \propto 1/T$ a $P$ constante cai à metade;
$\rho vS$ conserva-se: a celeridade dobra.
\end{solution}

\begin{solution}{exo:b2:fluid-kinematics:8}
(a) $\omega_z = \frac1r\frac{\dd(rv_\theta)}{\dd r}$: $2\Omega$ dentro, $0$
fora. (b) Dentro, $\Gamma = 2\pi r\cdot\Omega r = 2\Omega\cdot\pi r^2$, o fluxo
da \omterm{def:b2:fluid-kinematics:vorticity}{vorticidade}; fora, $\Gamma = 2\pi\Omega a^2$, constante, o fluxo de
todo o núcleo. (c) $\Omega = v_{\max}/a = \qty{1.4}{rad/s}$; $\Gamma = 2\pi av_{\max}
= \qty{2.2e4}{m^2/s}$; a \qty{500}{m}: $\Omega a^2/r = \qty{7}{m/s}$. (d) O
núcleo gira como um sólido; fora, o escoamento é irrotacional: as partículas
dão a volta ao eixo sem girar em torno de si mesmas. A \omterm{def:b2:fluid-kinematics:vorticity}{vorticidade} ---
a rotação --- mora no núcleo.
\end{solution}

\begin{solution}{exo:b2:fluid-kinematics:9}
(a) $\varphi = \tfrac12k(x^2 - y^2)$, $\Delta\varphi = k - k = 0$. (b) $\dd x/
kx = \dd y/(-ky)$: $xy = $ const.; $\partial_y\psi = kx = v_x$, $-\partial_x\psi =
-ky = v_y$. (c) $\vect a = (v_x\partial_x + v_y\partial_y)\vect v = k^2(x, y) =
\operatorname{\vect{grad}}\bigl(\tfrac12k^2(x^2 + y^2)\bigr) = \operatorname{\vect{grad}}
(v^2/2)$: num \omterm{def:b2:fluid-kinematics:vorticity}{escoamento irrotacional} permanente a aceleração é o
gradiente da energia cinética por unidade de massa; ela aponta para longe do
ponto de estagnação --- uma partícula desacelera ao se aproximar da parede
ao longo de $y$ e acelera ao se afastar ao longo de $x$. (d) $\dot x = kx$, $\dot y = -ky$:
$x = x_0\eu^{kt}$, $y = y_0\eu^{-kt}$ (com $xy$ constante).
\end{solution}

\begin{solution}{exo:b2:fluid-kinematics:10}
(a) $\dd x/\dd t = x/(t + \tau)$ dá $x = x_0(1 + t/\tau)$, $\ddot x = 0$;
pela \omterm{thm:b2:fluid-kinematics:material}{derivada material}: $\partial_tv + v\partial_xv = -x/(t + \tau)^2 + x/(t + \tau)^2
= 0$. (b) $\partial_t\rho + \partial_x(\rho v) = \dot\rho + \rho/(t + \tau) = 0$: $\rho =
\rho_0\tau/(t + \tau)$. (c) A separação de duas partículas cresce como $(1 + t/\tau)$
e a massa entre elas é fixa: $\rho \propto 1/(1 + t/\tau)$. (d) $v =
Hx$ com $H = 1/(t + \tau)$: uma lei de Hubble, expansão livre (sem
aceleração), com a massa específica caindo com o inverso do fator de escala.
\end{solution}

\begin{solution}{exo:b2:fluid-kinematics:11}
(a) $\vect v = \operatorname{\vect{grad}}\varphi = (U + qx/2\pi r^2, qy/2\pi r^2)$;
a fonte tem $v_r = q/2\pi r$, $\operatorname{div} = \frac1r\dd(rv_r)/\dd r
= 0$ e nenhuma dependência em $\theta$, logo $\omega_z = 0$; somas de campos assim são
de novo incompressíveis e irrotacionais. (b) No semieixo $x$ negativo,
$U - q/2\pi r = 0$: $x_s = -q/2\pi U$. (c) No ponto de estagnação $\theta =
\pi$, $y = 0$: $\psi = q/2$; a \omterm{def:b2:fluid-kinematics:euler}{linha de corrente} $Ur\sin\theta + q\theta/2\pi = q/2$,
ou seja, $r\sin\theta = (q/2\pi U)(\pi - \theta)$ --- uma curva que parte de $x_s$ e
se abre para jusante. (d) $\theta \to 0$: $y \to q/2U$; $\theta \to 2\pi$ (ramo
inferior): $y \to -q/2U$. O fluido dentro dessa \omterm{def:b2:fluid-kinematics:euler}{linha de corrente} vem da
fonte; fora dela, a corrente uniforme escoa em torno de um semicorpo de
largura assintótica $q/U$.
\end{solution}

\begin{solution}{exo:b2:fluid-kinematics:12}
(a) $\operatorname{div}\vect v = \frac1r\partial_r(-q/2\pi) = 0$; $\omega_z =
\frac1r\partial_r(\Gamma/2\pi) = 0$. (b) $\dd r/v_r = r\,\dd\theta/v_\theta$ dá
$\dd r/r = -(q/\Gamma)\dd\theta$. (c) $a_r = v_r\partial_rv_r - v_\theta^2/r =
-(q^2 + \Gamma^2)/4\pi^2r^3 = -v^2/r$; $a_\theta = v_r\partial_rv_\theta + v_rv_\theta/r
= 0$. (d) $\dd r/\dd t = -q/2\pi r$: $t = \pi(r_0^2 - a^2)/q$; voltas
$= (\Gamma/2\pi q)\ln(r_0/a)$. (e) $t = \qty{520}{s} \approx \qty{9}{min}$; $10$
voltas.
\end{solution}

\begin{solution}{pb:b2:fluid-kinematics:1}
\textbf{1.} Através de um cilindro de raio $r$: $2\pi rh|v_r| = D_V$, logo
$v_r = -q/2\pi r$, $q = D_V/h = \qty{1.5e-3}{m^2/s}$.

\textbf{2.} $rv_\theta = 0.5 \times 0.01 = \qty{5e-3}{m^2/s}$: $v_\theta = \Gamma/2\pi r$
com $\Gamma = \qty{3.1e-2}{m^2/s}$.

\textbf{3.} $\operatorname{div}\vect v = \frac1r\partial_r(rv_r) = 0$, $\omega_z =
\frac1r\partial_r(rv_\theta) = 0$.

\textbf{4.} $\dd r/r = -(q/\Gamma)\dd\theta$: $r = r_0\eu^{-q\theta/\Gamma}$; voltas
$= (\Gamma/2\pi q)\ln(r_0/a) = 3.3 \times 3.2 \approx 11$.

\textbf{5.} $t = \pi(r_0^2 - a^2)/q = \qty{520}{s} \approx \qty{9}{min}$; $v_\theta(a) =
\qty{0.25}{m/s}$.

\textbf{6.} $\vect a = -(v^2/r)\vect e_r$; em $r = a$: $v^2 = 0.25^2 + 0.012^2$,
$a = \qty{3.1}{m/s^2} \approx 0.3g$.

\textbf{7.} $\Omega_\oplus\sin\lambda\,r = 5 \times 10^{-5} \times 0.5 = \qty{2.5e-5}{m/s}$,
quatrocentas vezes menos que o movimento imposto: o hemisfério não decide
nada numa banheira comum (mas decide num tanque deixado quieto por um dia).

\textbf{8.} $v_\theta = \Omega r$ ($r < a$), $\Omega a^2/r$ ($r > a$); $\Omega = 80/60
= \qty{1.33}{rad/s}$, $\Gamma = 2\pi av_{\max} = \qty{3.0e4}{m^2/s}$.

\textbf{9.} $\omega_z = 2\Omega = \qty{2.7}{s^{-1}}$ no núcleo, $0$ fora;
$v_\theta$ cresce linearmente até \qty{80}{m/s} e depois cai como $1/r$; $\omega_z$ é um
degrau.

\textbf{10.} $a = v_{\max}^2/a = 6400/60 = \qty{107}{m/s^2} \approx 11g$, para dentro.

\textbf{11.} Fora, a \omterm{def:b2:fluid-kinematics:vorticity}{vorticidade} é nula: um elemento pequeno translada
ao longo de seu círculo sem girar em torno de si --- a rodinha de pás
mantém sua orientação. No núcleo (rotação sólida) ela gira com o
fluido, uma volta por revolução.

\textbf{12.} $E_{\text{nuc}} = \int_0^a\tfrac12\rho\Omega^2r^2\,2\pi r\,\dd r =
\tfrac14\pi\rho\Omega^2a^4 = \qty{2.2e7}{J/m}$.

\textbf{13.} $E_{\text{ext}} = \int_a^R\tfrac12\rho(\Omega a^2/r)^2\,2\pi r\,\dd r =
\pi\rho\Omega^2a^4\ln(R/a) = 8.7 \times 10^7 \times 4.4 = \qty{3.8e8}{J/m}$; o
integrando $\rho v^2r \propto 1/r$, donde o logaritmo.

\textbf{14.} $\qty{4.0e8}{J/m} \times \qty{1000}{m} = \qty{4e11}{J}$: cerca de
\qty{100}{t} de TNT.

\textbf{15.} Entrada $2\pi a \times 300 \times 5 = \qty{5.7e5}{m^3/s}$ através de
$\pi a^2 = \qty{1.1e4}{m^2}$: $v_z = \qty{50}{m/s}$.

\textbf{16.} $\omega = 2v_{\max}/a = \qty{2.5e-3}{s^{-1}} = 50f$; $\Gamma = 2\pi av_{\max}
= \qty{1.3e7}{m^2/s}$.

\textbf{17.} $r_0\cdot fr_0/2 = r(v_\theta + fr/2)$: $v_\theta = f(r_0^2 - r^2)/2r$;
\qty{60}{m/s} a \qty{100}{km}, \qty{155}{m/s} a \qty{40}{km}.

\textbf{18.} $f > 0$ no hemisfério norte dá $v_\theta > 0$:
sentido anti-horário visto de cima (ciclônico). Em \qty{500}{km} a
rotação planetária é a rotação inicial dominante; numa banheira é
o movimento imposto.

\textbf{19.} O atrito com o mar retira momento angular, e o
gradiente de pressão não consegue equilibrar a aceleração centrípeta de um
anel cada vez mais rápido: o ar sobe antes de alcançar o eixo, na parede
do olho, deixando um olho calmo.

\textbf{20.} $\rho\,2\pi rHv = 1.2 \times 2\pi \times 5 \times 10^5 \times 10^3 \times 5 =
\qty{1.9e10}{kg/s}$.

\textbf{21.} $\vect v(M, t) = -U\vect e_x + \vect v_0(M + Ut\,\vect e_x)$: o
padrão é advectado, e o campo não é permanente no referencial do solo.
A translação soma-se ao redemoinho do lado em que $v_\theta$ aponta para
oeste --- o lado norte (o da direita) de uma \omterm{def:b2:fluid-kinematics:euler}{trajetória} rumo a oeste.

\textbf{22.} $0.05 \times 10^7/3600 = \qty{139}{m^3/s}$; $v = 139/40 = \qty{3.5}{m/s}$.

\textbf{23.} $D_V = \qty{83}{cm^3/s}$: aorta \qty{0.33}{m/s}; capilares
\qty{0.33}{mm/s}; \qty{3}{s} por milímetro.

\textbf{24.} $S\,\dd h/\dd t = -s\sqrt{2gh}$: $\sqrt h = \sqrt{h_0} -
\frac sS\sqrt{g/2}\,t$, vazio em $T = (S/s)\sqrt{2h_0/g} = 10^4 \times 0.45 =
\qty{4500}{s} \approx \qty{75}{min}$.

\textbf{25.} Banheira: $r \sim \qty{0.5}{m}$, $v \sim \qty{1}{cm/s}$, $\Gamma =
\qty{3e-2}{m^2/s}$; tornado: \qty{60}{m}, \qty{80}{m/s}, $\omega = \qty{2.7}{s^{-1}}$,
$\Gamma = \qty{3e4}{m^2/s}$; ciclone: \qty{40}{km}, \qty{50}{m/s}, $\omega =
\qty{2.5e-3}{s^{-1}}$, $\Gamma = \qty{1.3e7}{m^2/s}$. Em cada um, o momento
angular por unidade de massa $rv_\theta$ do fluido que converge se conserva:
a convergência os acelera.
\end{solution}
